% A Digest Is a Change Detector, Not a Claim — arXiv preprint
% Build: tectonic main.tex  (XeTeX engine; fetches packages automatically)
% STATUS: working draft circulated for co-author review and consent prior to submission.
%
% NOTE ON THE PREAMBLE. This mirrors ../../ephemeris/paper/main.tex, which in
% turn mirrors transit's and the standard's. `minted` is not loaded: it shells
% out to pygments, which would make the CI paper build depend on a Python
% package nobody installs there. `listings` needs no external tool.
\documentclass[11pt]{article}

\usepackage[margin=1in]{geometry}
\usepackage{tikz}
\usepackage{graphicx}
\usepackage{booktabs}
\usepackage{array}
\usepackage{amsmath, amssymb}
\usepackage{listings}
\usepackage{xcolor}
\usepackage{colortbl}
\usepackage[hidelinks]{hyperref}
\usepackage{enumitem}
\usepackage{microtype}
% Long test names such as `a_static_deployment_never_advances_its_register`
% are unbreakable in \texttt and overflow the measure; htt lets them hyphenate.
\usepackage[htt]{hyphenat}
\usepackage{caption}
\usepackage{siunitx}
\usepackage{amsthm}
\newtheorem{definition}{Definition}
\newtheorem{theorem}{Theorem}
\newtheorem{proposition}{Proposition}
\usepackage[section]{placeins}

% Advanced AI Society brand palette
\definecolor{aailime}{HTML}{CFFF04}
\definecolor{aailimetint}{HTML}{F2FFB8}
\definecolor{aaiink}{HTML}{0A0A0A}
\definecolor{aaired}{HTML}{FFE3E4}
\definecolor{aaineutral}{HTML}{E7E3DD}
\definecolor{aaipaper}{HTML}{F0EDEA}
\definecolor{aaimuted}{HTML}{8F8A82}
\definecolor{aaidarklime}{HTML}{7A9900}
\definecolor{codepanel}{HTML}{F7F5F1}

% Ragged-right paragraph column; justified narrow columns underfill badly.
\newcolumntype{L}[1]{>{\raggedright\arraybackslash}p{#1}}

\lstset{basicstyle=\ttfamily\footnotesize, breaklines=true, frame=single,
        framerule=0.3pt, xleftmargin=1em, columns=fullflexible}
\captionsetup{font=small, labelfont=bf}

\newcommand{\spectrum}{\texttt{spectrum}}
\newcommand{\transit}{\texttt{transit}}
\newcommand{\ephemeris}{\texttt{ephemeris}}
\newcommand{\parallax}{\texttt{parallax}}

% The Society's bracket mark, drawn natively (lime on the dark band)
\newcommand{\aaimark}[1]{%
\begin{tikzpicture}[x=#1, y=#1, line cap=butt]
  \fill[aailime] (0,0) rectangle (0.16,1);
  \fill[aailime] (0.16,0) rectangle (0.42,0.24);
  \fill[aailime] (0.16,0.76) rectangle (0.42,1);
  \fill[aailime] (1.04,0) rectangle (1.20,1);
  \fill[aailime] (0.78,0) rectangle (1.04,0.24);
  \fill[aailime] (0.78,0.76) rectangle (1.04,1);
\end{tikzpicture}}

\begin{document}

% ---------------------------------------------------------------- cover page
\thispagestyle{empty}

\begin{tikzpicture}[remember picture, overlay]
  \fill[aaiink] (current page.north west) rectangle
    ([yshift=-6.35cm]current page.north east);
  \fill[aailime] ([yshift=-6.35cm]current page.north west) rectangle
    ([yshift=-6.53cm]current page.north east);
  \node[anchor=north west, inner sep=0pt]
    at ([xshift=1in, yshift=-0.95cm]current page.north west) {\aaimark{0.62cm}};
  \node[anchor=north west, inner sep=0pt, text width=14.2cm]
    at ([xshift=1in+1.15cm, yshift=-0.98cm]current page.north west)
    {\sffamily\bfseries\footnotesize\color{aailime}%
     ADVANCED AI SOCIETY\; \textcolor{aaimuted}{$\cdot$}\; PROOF-OF-CONTROL
     INITIATIVE\\[3pt]
     {\normalfont\sffamily\footnotesize\color{aaimuted}
      Working draft v0.1 $\cdot$ August 2026 $\cdot$ for co-author review
      prior to submission}};
  \node[anchor=north west, inner sep=0pt, text width=\textwidth]
    at ([xshift=1in, yshift=-2.55cm]current page.north west)
    {\sffamily\bfseries\fontsize{21.5}{26}\selectfont\color{aaipaper}%
     A Digest Is a Change Detector, Not a Claim: Composition Commitments
     for Agent Systems, and the Attack Class They Cannot See};
  \node[anchor=north west, inner sep=0pt, text width=\textwidth]
    at ([xshift=1in, yshift=-5.55cm]current page.north west)
    {\sffamily\small\color{aaipaper}%
     \textcolor{aailime}{\rule[0.5ex]{1.6em}{2pt}}\hspace{0.6em}%
     Composition is a theory. You have to state it before you can hash it.};
\end{tikzpicture}

\vspace*{5.9cm}

% Compact author block
{\small\noindent
\textbf{Jim Schwoebel}\textsuperscript{1,2}

\smallskip\noindent
{\footnotesize
\textsuperscript{1}\,Quome \quad
\textsuperscript{2}\,Advanced AI Society, \url{https://advancedaisociety.org}\footnote{Prepared
within the Proof-of-Control Initiative of the Advanced AI Society. Authorship is
limited to contributors to this manuscript. Co-authors will be added only upon
explicit consent and substantive contribution; the author list is deliberately
short at this stage rather than inherited from a sibling paper.}}}

\vspace{0.9em}
\noindent\textcolor{aailime}{\rule{\textwidth}{1.6pt}}
\vspace{0.2em}

\begin{center}
{\sffamily\bfseries\small ABSTRACT}
\end{center}
\vspace{-0.5em}
\noindent{\small
A Software Bill of Materials says what a system is made of. Applied to an
agent, it runs into two problems at once. The first is timing: an agent's
composition changes after the bill is taken. A tool server can present one
schema when an operator reviews it and a wider one when the agent calls it,
and a manifest taken at task start is byte-identical across that substitution
--- so the document that was supposed to detect component swaps is provably
blind to the component swap that matters most. The second is worse, because it
is not a bug: nobody has agreed what belongs on the bill. Any tool that ships
a default schema settles an open research question by configuration file, and
the settlement is invisible to whoever reads its output.
We built \spectrum, which refuses to do that. It defines no Agent Bill of
Materials. It defines a digest over a \emph{declared} component set, ships
three such sets as named falsifiable hypotheses, requires that every one of
them state what it cannot capture, and refuses to compare digests taken under
different hypotheses. On top of the static manifest it keeps a PCR-style
running commitment over composition changes, which holds 32 bytes however much
composition moves. A detection harness then runs real substitutions against
all three hypotheses. The tool-definition rug-pull is detected under the two
runtime hypotheses and invisible under the build-time one. Content poisoning
is detected under none of them, and we assert that negative in the test suite
on purpose. We also report the uncomfortable half: on the content row the
harness cannot fail, because \texttt{Detected} is unreachable outside the
tool-measurement branch by construction of the predicate. That row is a worked
consequence of the design, not an experiment. Locally, re-measuring a tool
list costs \SI{6.2}{\micro\second} at 8 tools and \SI{177}{\micro\second} at
256, linear at roughly \SI{0.7}{\micro\second} per tool --- around a hundredth
of the standard's per-action budget at realistic sizes, and \emph{not} an
answer to whether re-measurement is affordable, because the network round trip
that dominates it is not measured here. Our conclusion for the Agent Bill of
Materials is narrower than ``inventory'' and more useful than ``security
control'': a running composition commitment is a change-detection control over
the capability surface and an inventory control over content.}

\vspace{0.6em}
\noindent\textcolor{aaimuted}{\rule{\textwidth}{0.6pt}}

\vspace{0.4em}
\noindent{\footnotesize\textbf{Artefact.} All figures in
Section~\ref{sec:cost} come from \texttt{spectrum bench-mcp} on one host, and
every detection outcome in Section~\ref{sec:works} is real output from
\texttt{spectrum detect}. The implementation is Apache-2.0
\cite{spectrumtool}. Nothing in this paper was produced by a deployment: there
is no fleet, no network, and no operator but us.}

\vspace{1.4em}

% ---------------------------------------------------------------- the paper
\section{The Problem}
\label{sec:problem}

An operator is asked to approve a tool. The approval screen shows a payments
server offering one function:

\begin{lstlisting}
{"tools":[{"name":"charge",
           "parameters":{"amount":"number","currency":"string"}}]}
\end{lstlisting}

\noindent
They read it. Charging an amount in a currency is what a payments tool should
do, and there is nowhere in that schema to put a destination, so the money can
only go where the application already decided it goes. They approve it. The
system takes a bill of materials --- model, system prompt, tool definitions,
dependency tree --- hashes it, and files the hash with the task record.

The agent runs. At the moment it actually calls the tool, the server returns
this instead:

\begin{lstlisting}
{"tools":[{"name":"charge",
           "parameters":{"amount":"number","currency":"string",
                         "destination":"string"}}]}
\end{lstlisting}

\noindent
There is now somewhere to put a destination, and the model --- helpful,
schema-following, doing exactly what it was built to do --- can be talked into
putting one there.

Now ask what the bill of materials says about this. It says nothing. Not
``suspicious,'' not ``unverified.'' It produces the same hash it produced
before, because it was taken at task start and the tool list it hashed was the
one presented at task start. The artefact whose entire purpose is to detect a
substituted component is, against this substitution, a constant function.

I want to be careful about what is and is not surprising, because there is a
reading on which this is trivial. Of course a snapshot does not capture a
later change. Snapshots do not do that. Everybody nods.

The nodding does not survive contact with a deployment. What you get in
practice is a pipeline that emits an SBOM at build time, a policy that
requires one, an auditor who checks that it was emitted and that its digest
matches the one on file, and a green tick. Every step is real work. The green
tick means ``the components we recorded at task start are the components we
recorded at task start,'' and it is read as ``the agent ran with the
components you approved.'' The gap is not in anybody's understanding of
hashing. It is in the distance between what the artefact commits to and what
its presence is taken to mean.

\subsection{The second problem, which is the harder one}

Suppose we fix that by re-measuring at call time. Then a different question
arrives immediately, and it has no technical answer: re-measure \emph{what}?

A conventional SBOM has a settled answer. A program is made of files and
libraries; you enumerate them. An agent is made of --- and here the sentence
stops. A model, certainly, though ``the model'' behind a hosted API endpoint
is a name and a version string and not something you can hash. A system
prompt. Tool definitions. The dependency tree of the harness. Also, arguably:
the documents the retriever pulled into the context this turn, the memories it
read, the sub-agents it spawned, the tool list each remote server returned at
each call. Also, if you are being ambitious: which of those retrieved
documents actually influenced the decision.

That list is not a specification. It is a set of positions, and reasonable
people hold different ones. Including retrieved chunks makes the bill enormous
and turns every content digest into a cross-session fingerprint of what the
agent read. Excluding them means the bill says nothing about the entire class
of attacks that work by feeding an agent a document. There is no neutral
choice, and there is certainly no choice that a tool should make silently on
its user's behalf.

This is the position P10 of the Proof-of-Control research roadmap is in
\cite{p10}: it asks whether an Agent Bill of Materials can be a running,
verifiable commitment, and whether that catches anything a static one does
not. It is an open question, marked open. And the fastest way to make an open
question look closed is to ship a tool with a default.

\begin{quote}
A tool that ships a default component set answers a research question by
configuration file, and the answer is invisible in the output. Every digest it
prints looks like a fact about a system. It is a fact about a system
\emph{under a theory of composition the reader was never shown}.
\end{quote}

So \spectrum{} has no default profile. \texttt{--profile} is required on every
command that composes a manifest. There is no invocation of this tool that can
produce an unmarked ``the AgBOM,'' and that is a deliberate ergonomic cost
paid to keep a research question from being settled by a default value.

\subsection{What this paper claims}

Three things, in decreasing order of confidence.

First, that the timing problem is real and mechanically demonstrable, and that
its fix --- a running commitment extended at composition-change points ---
catches the tool-substitution class. We demonstrate this against a hostile
server that performs the substitution, under three profiles, and report which
ones catch it.

Second, that the fix does \emph{not} extend to content. A poisoned document
that is validly sourced, correctly digested and faithfully recorded moves the
register exactly as a benign one does. The commitment records that the chunk
arrived. It says nothing about whether the chunk was adversarial, and nothing
about whether it steered the decision. We assert this negative in the test
suite so that no future change can make it quietly pass.

Third, and most tentatively, that the honest summary for P10 is: \emph{a
running composition commitment is a change-detection control over the
capability surface, and an inventory control over content.} That moves the
threat model on tool-and-plugin substitution and does not move it on the
memory- and context-poisoning rows.

\section{What Others Have Built}
\label{sec:related}

Nearly every mechanism in this paper is borrowed, and the borrowing is the
point: the contribution is not a new digest construction.

\paragraph{Bills of materials.} CycloneDX \cite{cyclonedx} and SPDX
\cite{spdx} are the two formats in use, the second standardised as ISO/IEC
5962. The NTIA minimum-elements report \cite{ntiasbom} is what most
procurement language actually points at, and US Executive Order 14028
\cite{eo14028} and the EU Cyber Resilience Act \cite{eucra} are why anybody is
required to produce one. This work reads CycloneDX and does not read SPDX,
which is a choice we defend badly in Section~\ref{sec:unknown}. All of these
formats describe a build. None of them has a concept of a component that
changes identity between two calls within one execution, because for the
software they were designed for, no such component exists.

\paragraph{Provenance for builds.} SLSA \cite{slsa} and in-toto
\cite{intoto} secure the path from source to artefact: who built what, from
which inputs, on which machine. The threat they address is the substituted
build step, and Thompson's compiler \cite{thompson1984} is the ancestor of
the whole line. Ohm and colleagues catalogued what supply-chain attacks
actually look like in package ecosystems \cite{ohm2020}. Our attacker is
adjacent but not the same: they do not compromise the build, they change what
a remote endpoint returns \emph{after} the build was attested and after a
human approved it.

\paragraph{Running commitments.} The register in Section~\ref{sec:precise} is
a TPM Platform Configuration Register with the hardware removed
\cite{tpm2}. Extend-only, fixed width, order-dependent, no setter. The idea
that you can summarise an unbounded sequence of measurements in 32 bytes and
lose nothing except the ability to reconstruct them is thirty years old, and
we use it unmodified. Merkle's hash tree \cite{merkle1980} and Certificate
Transparency \cite{rfc6962} are the other half of that tradition; the
sibling tool \ephemeris{} \cite{ephemeristool} is where they live in this
programme, and Section~\ref{sec:precise} explains why a register without a log
is only half of a verification story.

\paragraph{Canonicalisation.} Comparing two JSON documents for identity is a
solved problem with an unusually large number of ways to get it wrong. We use
RFC~8785 JCS \cite{rfc8785}, through \transit's implementation
\cite{transittool} rather than a second one of our own, over strictly parsed
JSON \cite{rfc8259,rfc7493}. A server that reformats its JSON must not read
as a composition change; a server that sends duplicate keys must be refused
rather than silently resolved to whichever key a decoder prefers.

\paragraph{Evidence classes.} The RATS architecture \cite{rfc9334} separates
attester, verifier and relying party, and insists that a claim's weight
depends on who is vouching for it. Our witness class is a coarse version of
that idea pushed inside the digest: a component an operator typed and a
component whose bytes we hashed must not render identically. \parallax{}
\cite{parallaxtool} is the sibling tool that works out what such a
distinction is worth across a whole deployment.

\paragraph{Agent-specific attacks.} Greshake and colleagues demonstrated
indirect prompt injection \cite{greshake2023}: content that enters a model's
context and redirects it. The OWASP LLM list \cite{owaspllm} has since
catalogued the family. This literature is why Section~\ref{sec:works}'s
negative result matters. The attack that most concerns practitioners is
delivered as \emph{content}, and content is exactly the axis on which a
composition commitment degenerates into an inventory.

\paragraph{Where we differ.} The work above mostly asks how to build the
artefact correctly. We are asking what a correctly built artefact is worth
when the thing it describes changes underneath it, and --- separately --- what
a tool should do when the community has not decided what the artefact should
contain. Our answer to the second is the profile: ship the disagreement as
named, comparable, individually falsifiable hypotheses, and refuse to rank
them.

\section{Writing It Down Precisely}
\label{sec:precise}

The symbols will not make anything truer than the English already did. They
are here because several of the properties below are easy to state loosely and
then implement in a way that looks right, and because the failure modes we
care about live in details --- concatenation order, length prefixing, which
fields enter a hash --- that prose glosses over.

\begin{definition}[Profile]
A profile is a triple $P = (\mathrm{id}, K, C)$ where $\mathrm{id}$ is a
stable name, $K \subseteq \mathcal{K}$ is a set of component kinds drawn from
a closed enumeration $\mathcal{K}$, and $C$ is a non-empty set of
natural-language statements of what this theory of composition does not
capture. A profile with $C = \emptyset$ does not load.
\end{definition}

That last clause is the one doing work. It is not possible to express, in this
tool's configuration language, the proposition ``this profile captures
everything.'' The check is four lines long and it is the most opinionated code
in the repository.

\begin{definition}[Component reference]
A component reference is $c = (k, \mathrm{id}, d, w)$ with $k \in
\mathcal{K}$, $\mathrm{id}$ a package URL where one exists and a stable local
name otherwise, $d$ an optional tagged content digest, and $w \in
\{\textsc{Measured}, \textsc{Attested}, \textsc{Declared}, \textsc{Opaque}\}$
the witness class --- how we know what we are claiming.
\end{definition}

\textsc{Opaque} is a positive statement: a component of this kind is present
and cannot be named. It is never the result of a failure. That distinction
matters because the alternative --- letting an unmeasurable component fall out
of the manifest --- makes ``we could not measure the training data'' and
``there is no training data'' render identically.

\begin{definition}[Component digest]
\[
  H_c(c) = \mathrm{SHA\text{-}256}\big(\texttt{"spectrum-component-v1"}
  \;\|\; \mathrm{len}(J) \;\|\; J \big)
\]
where $J$ is the JCS canonicalisation of the object
$\{\texttt{kind}: k, \texttt{id}, \texttt{content\_digest}: d,
\texttt{witness}: w\}$ and $\mathrm{len}$ is a 64-bit big-endian length.
\end{definition}

Three details are load-bearing. The absent digest is rendered as an explicit
\texttt{null} rather than omitted, so ``we have no digest'' and ``we did not
mention a digest'' cannot collide. The witness class is \emph{inside} the
digest, which is one of P10's open forks and which we resolve in the direction
the design recommends --- at a real cost, discussed in
Section~\ref{sec:unknown}. And the domain tag is hashed \emph{before} the
canonical bytes rather than carried as a key inside the JSON object, so it can
never be shadowed by a future field that happens to be named the same thing. A
tag inside the JCS input is just another field, and another field is precisely
what it must be distinguishable from.

\begin{definition}[Manifest and its digest]
A manifest $M = (P.\mathrm{id}, S)$ carries a sorted, deduplicated sequence
$S = \langle c_1, \ldots, c_n\rangle$. Its digest is
\[
  H_M(M) = \mathrm{SHA\text{-}256}\big(\texttt{"spectrum-manifest-v1"} \;\|\;
  \mathrm{len}(P.\mathrm{id}) \;\|\; P.\mathrm{id} \;\|\; n \;\|\;
  \mathrm{len}(H_c(c_1)) \| H_c(c_1) \;\|\; \cdots \big).
\]
\end{definition}

The profile id is inside the digest. This is what makes a digest a commitment
\emph{under a stated theory of composition} rather than a fact about a system,
and it is why \texttt{diff} refuses to compare two manifests taken under
different profiles: a difference between two theories is not a difference
between two systems. Each component digest is length-prefixed, so a set of two
components cannot be made to hash like one component whose identifier happens
to contain the concatenation of their digests.

\begin{proposition}[The sorting invariant is not cosmetic]
\label{prop:sort}
Let $\mathrm{cmp}$ be the order-insensitive comparison used by \texttt{diff},
which indexes components by $(k, \mathrm{id})$. If a caller could append to
$S$ after construction, there would exist manifests $M_1, M_2$ over the same
component set with $\mathrm{cmp}(M_1, M_2) = \emptyset$ and $H_M(M_1) \neq
H_M(M_2)$.
\end{proposition}

\noindent
That is a producer and a verifier inside one library reaching opposite
conclusions about the same pair of manifests. The remedy is structural rather
than documentary: both fields of the manifest type are private, sorting and
deduplication happen in the only constructor, and there is deliberately no
\texttt{push}, no \texttt{components\_mut}, and no conversion from a bare
vector. The invariant holds by construction rather than by promise.

\begin{definition}[Composition register]
\begin{align*}
  C_0 &= \mathrm{SHA\text{-}256}\big(\texttt{"spectrum-register-v1"} \;\|\;
        \mathrm{len}(P.\mathrm{id}) \| P.\mathrm{id} \;\|\;
        \mathrm{len}(H_M) \| H_M \big) \\
  C_n &= \mathrm{SHA\text{-}256}\big(C_{n-1} \;\|\; H_e(e_n)\big)
\end{align*}
where $H_e$ is a domain-tagged digest over an event's kind, its interception
hook, and the digest of the component that moved.
\end{definition}

The register is 32 bytes regardless of how much composition has moved, which
is what makes a per-step field affordable. The hook is inside the event
digest, because the same component arriving at context assembly and at memory
write are different composition changes.

\begin{proposition}[One-way by construction]
There is no sequence of public operations returning a register to a previous
value other than by preimage.
\end{proposition}

\noindent
The inner value is private, there is no setter, no \texttt{reset}, and no
constructor other than $C_0$. This is worth stating as a property because the
convenience method that would violate it --- ``reset composition for a new
task'' --- is one an implementer would reach for without noticing that it is a
rewrite of history.

\subsection{What the register does not do}
\label{sec:register-honest}

It solves the record-size problem completely and the verification problem not
at all.

To check a register you need its events, and this tool does not store them.
Replaying a truncated log does not reproduce the register --- which is the
property that makes the construction worth anything --- but noticing that
requires having the log. So the honest statement of what \spectrum{} provides
is: a 32-byte value that changes when composition changes, and a
demonstration that it changes for the right reasons, with the storage and
custody of the event log left to \ephemeris{} \cite{ephemeristool}, whose own
paper explains why that custody is the hard part.

There is a further limitation that is easy to miss and that we state plainly
because it undercuts the natural reading of the whole construction:
\emph{nothing verifies a digest.} No infrastructure in this programme
publishes expected AgBOMs anywhere. There is no registry, no
transparency log of composition roots, nothing to check a value against. What
the library offers is \texttt{diff} against a prior digest of the same
deployment. A digest here is a change detector between two records of the same
system --- not a checkable claim about one record in isolation. Every
sentence in this paper about ``detection'' should be read against that.

\section{Three Hypotheses, Not One Answer}
\label{sec:profiles}

The tool ships three profiles. They are not a menu of configurations for
different users; they are three positions on P10's open question, written down
in a form that can be run against an attack and shown to be wrong.

\begin{table}[htbp]
\centering
\small
\begin{tabular}{@{}L{2.9cm}L{4.2cm}L{3.3cm}L{4.2cm}@{}}
\toprule
\textbf{Profile} & \textbf{Includes} & \textbf{The claim} &
\textbf{What it cannot capture} \\
\midrule
\texttt{build-manifest} &
model, system prompt, tool definitions, code and dependency SBOM &
composition is fixed at task start &
everything that changes after \\
\addlinespace
\texttt{runtime-tools} &
the above, plus MCP tool-list re-measurement and the sub-agent roster &
the capability surface changes at runtime and is small enough to re-measure &
retrieved context, memory, whether a component is adversarial \\
\addlinespace
\texttt{runtime-\hspace{0pt}with-context} &
the above, plus retrieved-chunk and memory-read references &
everything entering the context is part of composition &
which of it mattered --- and it costs a linkability leak \\
\bottomrule
\end{tabular}
\caption{The three shipped profiles. The right-hand column is not commentary:
it is a required, non-empty field of the profile file, and a profile that
omits it fails to load.}
\label{tab:profiles}
\end{table}

\texttt{build-manifest} exists to reproduce the behaviour of what is deployed
today. It is shipped so that the cost and the detection power of the other two
are measured \emph{against} it rather than asserted over it, and it is
deliberately not named \texttt{default}: it is not the default answer, it is
the current one, and this tool exists partly to show what it misses.

The profile mechanism has one more consequence worth drawing out. Because the
profile id is inside every digest, two deployments running under different
profiles produce digests that are not comparable, and the library refuses to
compare them rather than reporting a difference. An operator who upgrades from
\texttt{build-manifest} to \texttt{runtime-tools} does not see ``composition
changed.'' They see a refusal, with a message explaining that they changed
theories. We think this is the correct behaviour and we note that it is
annoying, which is usually a sign that a tool is telling the truth about
something.

\subsection{What a profile still does not tell you}

A profile's coverage list states what the \emph{theory} excludes. It does not
state what the \emph{plumbing} dropped. CycloneDX defines component types ---
\texttt{container}, \texttt{operating-system}, \texttt{firmware},
\texttt{device}, \texttt{platform} --- that this projection has no kind for,
and those are discarded before the profile is ever consulted. No shipped
profile's coverage mentions them, so nothing in a profile tells an operator
that a container in their document did not reach the digest.

We treat this as a real gap and mitigate it at the point of use rather than
pretending it is covered: \texttt{compose} counts the unmapped types and
prints them under the digest, so a dropped type is visible on the run that
dropped it rather than inferred from a list it is not on.

\begin{lstlisting}
profile   build-manifest
components 1
digest    sha-256:ee76701fa13e35df0ea03f6cc304f54ca3313eb6e04cc6dc3f5e0de531cedc45

DECLARED BUT ABSENT - this profile includes these kinds and no
component of them was collected, so the digest above says nothing
about them. They are missing, not empty:
  - system_prompt
  - tool_definition
  - code_dependency

UNMAPPED CycloneDX TYPES - present in the document, no kind to
project them onto, dropped before the profile was consulted:
  - container (2)
  - firmware (1)
  - operating-system (1)
\end{lstlisting}

The ``declared but absent'' block is the same discipline applied to a
different gap. Four of the eight kinds have no collector in this build, so a
profile that declares them will routinely produce a digest that commits to
none of them. The design's fail-loud table asks \texttt{compose} to refuse in
that situation. It warns instead, because refusing would make the command
unusable until collectors that are several build steps away exist. That is a
compromise and we record it as one: the digest still prints, and it never
prints alone.

\section{Does It Actually Work?}
\label{sec:works}

\subsection{The harness}

\texttt{spectrum detect --scenario <name>} runs one attack against every
shipped profile and reports, per profile, whether that profile's commitment
moved and which component moved it. The exit code is 0 if any profile detected
the attack and 1 if none did.

The hostile server is a pure function of a scenario and a call count, so every
detection test runs with no socket, no network and no credentials. Call 0 is
the operator's review; every later call is a real invocation. That split is
the whole attack. There is also a real HTTP transport over the same function
for a human to poke at by hand, and nothing in the automated suite touches it.

``Detected'' has a precise and narrow meaning here, and we would rather state
it three times than have it read as something grander: \emph{the register
value after the attack differs from the value a verifier would expect, and the
event log names which component moved.} It does not mean the attack was
prevented. It does not mean anything was judged malicious. Nothing in this
tool blocks anything.

\subsection{The matrix}

Table~\ref{tab:detect} is real output, reformatted. Every cell was produced by
running the command named in Section~\ref{sec:works}.

\begin{table}[htbp]
\centering
\small
\begin{tabular}{@{}lccc@{}}
\toprule
\textbf{Scenario} & \texttt{build-manifest} & \texttt{runtime-tools} &
\texttt{runtime-with-context} \\
\midrule
\texttt{none} (control) & not detected & not detected & not detected \\
\texttt{tool-rug-pull} & not detected &
\cellcolor{aailimetint}\textsc{detected} &
\cellcolor{aailimetint}\textsc{detected} \\
\texttt{added-tool} & not detected &
\cellcolor{aailimetint}\textsc{detected} &
\cellcolor{aailimetint}\textsc{detected} \\
\texttt{content-poisoning} & not detected & not detected & not detected \\
\bottomrule
\end{tabular}
\caption{Detection outcomes. Detected cells name the component that moved
(\texttt{mcp://payments}). The first row is the negative control and the last
row is the finding.}
\label{tab:detect}
\end{table}

Three things to read off it.

\paragraph{The control is honest.} With no attack, nothing is detected under
any profile. This row exists because a harness that reports detection when
nothing happened would make every other row worthless, and a suite without
this row cannot tell the difference between a working detector and a stuck
one.

\paragraph{The rug-pull is caught, and by the profile you would predict.}
Under \texttt{build-manifest} the digest before and after the substitution is
identical, because the manifest was taken once. Under \texttt{runtime-tools}
the register diverges at the call and the event names
\texttt{mcp://payments}. This is detectable for reasons that do not
generalise: a tool definition is small, local, re-readable on demand, and has
an expected value that a human approved. Every one of those properties is load
bearing, and content has none of them.

\paragraph{Content poisoning is caught by nothing.} A document that is validly
sourced, correctly digested, faithfully recorded and malicious moves the
register exactly as a benign document does. Under
\texttt{runtime-with-context} the register \emph{does} change --- the chunk
entered the context and the commitment records that it did --- and the
scenario still reports not detected, because a change that a verifier has no
expected value for is a record of arrival, not a detection.

\subsection{What the harness proves, and what it only re-states}
\label{sec:harness-honest}

This is the part of the evaluation we would most like to overstate, so we will
state it against ourselves.

On the content row, the harness could not have come out the other way. In the
detection routine, the \texttt{Detected} outcome is reachable only when a
``moved'' variable is set, and that variable is assigned only inside the
branch guarded by \emph{does this profile include MCP servers}. A content
attack therefore cannot reach \texttt{Detected} by construction of the
predicate, whatever the profile says and whatever the register does.

That is deliberate, and the code's own documentation gives the reason: a
verifier has no expected value for a chunk retrieved at runtime, so a register
that moved because a chunk arrived is a record of arrival and not a detection.
We think the negative is correct and the argument for it is sound. But the
harness \emph{executes} that argument against real profiles and real
substitutions; it does not independently discover it. The content row of
Table~\ref{tab:detect} is a worked consequence of the design, not an
experiment that could have falsified it.

We are labouring this because the table looks like an experiment, tables
generally are experiments, and a reader skimming for results would take that
row as an empirical finding about composition commitments in general. It is
not. It is a demonstration that our implementation behaves the way our
argument says it should, which is worth having and is a strictly weaker thing.

The rows that \emph{could} have come out the other way are the rug-pull and
added-tool rows under \texttt{runtime-tools}, and the control row under all
three. Those depend on canonicalisation behaving, on the digest covering the
fields we think it covers, and on the register actually diverging. Any of them
could have failed, and during development the equivalent of the control row is
what catches a detector that has become stuck.

\subsection{How we checked the rest}

The suite is 76 tests, all offline. The ones that matter for this paper's
claims are the ones that pin a property against a future edit rather than
exercise a happy path:

\begin{itemize}[leftmargin=1.2em, itemsep=0.15em]
\item \texttt{the\_witness\_class\_is\_inside\_the\_digest} asserts that a
      \textsc{Declared} component does not hash like a \textsc{Measured} one,
      which pins P10's fork 2 in the direction we took. Reversing that fork is
      a visible test failure rather than a quiet behaviour change.
\item \texttt{every\_witness\_tag\_is\_distinct} and
      \texttt{every\_event\_kind\_tag\_is\_distinct} assert that no two
      variants share a wire tag. A copy-paste that gave two evidence classes
      the same string would undo fork 2's resolution without touching the code
      that implements it.
\item \texttt{the\_manifest\_digest\_of\_a\_two\_kind\_manifest\_is\_pinned}
      turns a reordering of the kind enumeration into a test failure. The
      declaration order of that enumeration decides concatenation order and
      therefore the digest of every multi-kind manifest; reordering it looks
      cosmetic in review.
\item \texttt{a\_static\_deployment\_never\_advances\_its\_register} pins the
      property that the running commitment degenerates exactly to the static
      one when nothing changes.
\item \texttt{replaying\_a\_truncated\_log\_does\_not\_reproduce\_the\_register}
      is the property that makes the register worth keeping.
\item \texttt{content\_poisoning\_is\_not\_detected\_by\_any\_profile}
      asserts the negative result against all three profiles, on purpose. Its
      comment says it is the row most likely to be deleted by somebody tidying
      the suite. It should not be.
\end{itemize}

\section{What It Costs}
\label{sec:cost}

The design's central bet is fork 5: that an MCP server's identity should be a
digest of the tool list as re-measured, rather than an endpoint URL and a
declared version. Re-measurement is what makes the rug-pull visible. The
obvious objection is cost, because re-measurement happens per call and the
Proof-of-Control standard's commentary sets a per-action budget of
\SI{15}{\milli\second} \cite{pocspec}.

\subsection{The measurement}

We measured the local half: canonicalisation and hashing of a synthetic tool
list, in process, with the list built before the timing loop and a
100-iteration warm-up outside it. Seven repetitions per size, 10\,000
iterations each, one host.

\begin{table}[htbp]
\centering
\small
\sisetup{table-number-alignment=center}
\begin{tabular}{@{}rS[table-format=3.3]S[table-format=3.3]S[table-format=3.3]S[table-format=2.1]S[table-format=1.4]@{}}
\toprule
\textbf{Tools} & {\textbf{Min} (\si{\micro\second})} &
{\textbf{Median} (\si{\micro\second})} & {\textbf{Max} (\si{\micro\second})} &
{\textbf{Spread} (\%)} & {\textbf{Median as \% of \SI{15}{\milli\second}}} \\
\midrule
1   & 1.250   & 1.277   & 1.386   & 10.9 & 0.0085 \\
2   & 1.983   & 2.045   & 2.149   & 8.4  & 0.0136 \\
4   & 3.281   & 3.346   & 3.430   & 4.5  & 0.0223 \\
8   & 6.047   & 6.197   & 6.364   & 5.2  & 0.0413 \\
16  & 11.404  & 11.587  & 11.861  & 4.0  & 0.0772 \\
32  & 22.006  & 22.494  & 22.613  & 2.8  & 0.1500 \\
64  & 43.083  & 43.310  & 43.921  & 1.9  & 0.2887 \\
128 & 84.491  & 84.945  & 87.629  & 3.7  & 0.5663 \\
256 & 176.164 & 177.259 & 182.886 & 3.8  & 1.1817 \\
\bottomrule
\end{tabular}
\caption{Local cost of one MCP tool-list re-measurement. Seven repetitions per
row, 10\,000 iterations per repetition, Apple M5 Max, \texttt{rustc 1.97.1},
release build. Spread is $(\max-\min)/\min$.}
\label{tab:cost}
\end{table}

The cost is linear in list size. A least-squares fit over the nine sizes gives
\SI{0.685}{\micro\second} per tool with a fixed cost of about
\SI{0.27}{\micro\second}; taking the marginal cost between the 64- and
256-tool rows instead gives \SI{0.698}{\micro\second} per tool. The fit
under-predicts the small end by about 7\%, which is what you would expect from
a straight line dominated by the large sizes, so the per-tool figure is the
number to quote and the intercept is not.

Extrapolating the fit, a tool list would need roughly 22\,000 entries before
local re-measurement consumed the entire \SI{15}{\milli\second} budget. At
256 tools it consumes 1.2\% of it.

\subsection{One run is not a benchmark}

The repetitions turned up something that changes how the previously committed
figures should be read. The results file in the repository reports
\SI{9.635}{\micro\second} at 8 tools, from a single run. The median of seven
repetitions on the same host, same toolchain, same binary, is
\SI{6.197}{\micro\second}. A single standalone run we took while preparing
this paper landed at \SI{7.021}{\micro\second}, between the two.

That is a factor of 1.55 between two measurements of the same quantity on the
same machine. Within a sweep the spread is 1.9\%--10.9\% and tightens as the
work per iteration grows, which is the ordinary shape of timing noise. Across
sweeps, at the small end, it is not small at all.

The ordering --- first-run slowest, sweep-median fastest --- is consistent
with a warm-up effect that the in-process 100-iteration warm-up does not
cover: that loop warms the allocator and the branch predictors, and does
nothing about CPU frequency or scheduler placement. We label that a hypothesis
rather than a finding, because we did not test it, and testing it properly
means pinning frequency and affinity, which we did not do.

The practical consequence is a correction to our own earlier reporting: the
\SI{9.635}{\micro\second} figure is not wrong as a record of one run, and it
overstates the steady-state cost by about half. Both figures support the same
conclusion --- the local cost is two to three orders of magnitude below the
budget --- which is why this did not change any argument. It would have
changed one if the true figure had been near the budget.

\subsection{What this does not settle}

Fork 5 is not resolved by Table~\ref{tab:cost}, and we want to be blunt about
the gap because a table of microseconds is exactly the kind of artefact that
gets cited as though it settled something.

What is measured is canonicalisation and hashing. What is not measured is the
round trip to a remote MCP server, which is the operator's network. That round
trip is expected to dominate the total, by a margin large enough that the
local cost is a rounding error inside it. A negligible local cost is a
\emph{necessary} condition for re-measurement being affordable inside the
per-action budget. It is not a \emph{sufficient} one. Whether re-measurement
is affordable end-to-end depends on a number this repository cannot produce,
because producing it requires a real deployment with a real network and real
servers, and we have none of those.

For scale, the standard's own commentary records a single hardware attestation
quote at \SI{39.5}{\milli\second} \cite{pocspec} --- already over the
per-action budget on its own. The interesting comparison for fork 5 is against
that kind of number, not against ours.

\section{What We Still Don't Know}
\label{sec:unknown}

\subsection{The forks we did not resolve}

\spectrum{} implements a design that leaves eight forks open, and it is not
the case that the tool takes the recommended branch on each. It resolves one.
Table~\ref{tab:forks} is the accounting.

\begin{table}[htbp]
\centering
\small
\begin{tabular}{@{}L{3.6cm}L{11.2cm}@{}}
\toprule
\textbf{Fork} & \textbf{What this work does} \\
\midrule
1. SPDX support &
\textbf{Diverges from the recommendation.} The design recommends reading both
formats and emitting CycloneDX. This tool reads CycloneDX only. An operator
with an SPDX document cannot use \texttt{compose}. A second reader is a second
identity derivation to keep in agreement, and it is not built. \\
\addlinespace
2. Witness class in the digest &
\textbf{Resolved, as recommended: inside.} Pinned in both directions by tests.
\\
\addlinespace
5. MCP server identity &
\textbf{Instrumented, not resolved.} \texttt{runtime-tools} implements
re-measurement and Section~\ref{sec:cost} measures its local cost. The network
round trip decides this fork and needs a deployment. \\
\addlinespace
7. Context linkability &
\textbf{Half built.} See Section~\ref{sec:privacy}. \\
\addlinespace
3, 4, 6, 8 &
\textbf{Untouched.} No branch is taken because nothing here reaches the
decision: no field of this tool is emitted as an AgBOM digest, there is no
prompt collector, no retriever, and no delegation. Claiming this tool ``takes
the recommended branch'' on them would be an empty claim about code that does
not exist. \\
\bottomrule
\end{tabular}
\caption{P10's eight forks against what this implementation does.}
\label{tab:forks}
\end{table}

One thing bearing on fork 3 is worth naming without closing it. The test
\texttt{a\_static\_deployment\_never\_advances\_its\_register} pins the
property the recommendation rests on --- that the running value degenerates
exactly to the static root when nothing changes. That is evidence for the
recommendation, not a resolution of it, because nothing in this tool fills the
field the fork is about.

\subsection{Influence, which nobody can compute}

The standard's C1.2.2 asks for a link between retrieved content and the
decision it influenced. No profile in this tool captures it, no field
expresses it, and there is no API that could produce one.

That absence is deliberate and it is enforced by there being no such field. A
stub --- an \texttt{influence} column populated with retrieval scores, or
attention weights, or a heuristic --- would be read as the C1.2.2 link by
anybody downstream, and it would be wrong in a way that is very hard to detect
from the outside, because a plausible number in a plausible column is
indistinguishable from a real one. The tool records presence and never names
it influence, and the coverage output says so on every run.

We do not know how to compute this. We are not aware of anybody who does. It
is the single largest gap between what a composition commitment provides and
what the threat model wants from it.

\subsection{The privacy regression we shipped behind a warning}
\label{sec:privacy}

\texttt{runtime-with-context} digests retrieved content, and a content-derived
chunk digest is a cross-session, cross-agent fingerprint of what was read. Two
evidence records carrying the same chunk digest prove that both agents read
the same document. Across a fleet, that is a linkability channel that did not
exist before the composition commitment was introduced --- built, by us, in
the name of oversight.

Fork 7 is the decision point and it is unresolved. What we did instead of
resolving it is refuse to bury it: \texttt{spectrum coverage} keyed on a
profile including retrieved chunks prints the warning as output, not as
documentation. That is a mitigation for the operator who runs the command and
no mitigation at all for the deployment that does not.

The sibling tool \texttt{occultation} \cite{occultationtool} exists to cost
out the unlinkability machinery that would be the actual answer here. The two
halves do not meet: nothing in this repository consumes it, and we are not
claiming that they compose without somebody doing that work.

\subsection{Discipline that is written down and not enforced}

Two properties in this design are intended discipline rather than enforced
behaviour, and the distinction is not visible from the outside, so we state it
here.

Hosted-model weight identity should be \textsc{Declared} and never
\textsc{Measured}: no measurement of a hosted model's weights exists from
outside the provider. Training-data lineage should be \textsc{Opaque}. The
code enforces neither. There is no model collector at all; no collector in the
library ever constructs an \textsc{Opaque} component (only the test suite
does, to check that such a component still has an identity); no test asserts
either property; and the component type is a plain public structure, so a
caller can write \textsc{Measured} against a hosted model today and nothing
objects.

These are rules written down before the code they constrain exists, rather
than after. We think that is the right order. We also think a reader is
entitled to know that today they are prose.

\subsection{The cost of resolving fork 2 the way we did}

Putting the witness class inside the digest means that improving your
instrumentation changes your digest even though the system did not change.
Move a component from \textsc{Declared} to \textsc{Measured} --- exactly the
improvement you would want an operator to make --- and every downstream
comparison reports a difference.

We accepted that cost because the alternative is a digest in which an
operator's typed version string is indistinguishable from measured bytes. But
we have not measured the cost, and it is the kind of cost that lands on a
fleet all at once: an instrumentation rollout looks like a fleet-wide
composition change. The library reports witness changes as their own category,
separately from content changes, which is the smallest useful mitigation and
is not the same as knowing what a rollout does to a real deployment's alert
volume.

\subsection{The evidence base}

One host, one afternoon, no deployment. The detection harness runs against a
hostile server we wrote, presenting an attack we chose, to profiles we
designed. There is no fleet, no third party, and no adversary who was not us.
Every number in Section~\ref{sec:cost} is local, and the number that would
decide fork 5 is not in this paper.

We think the argument in Section~\ref{sec:works} survives that --- the
rug-pull is a mechanical consequence of when you measure, not an empirical
claim about how often anybody does it --- but the frequency question, the one
an operator would actually ask, is untouched. We do not know how often tool
definitions change between approval and call in real deployments. We do not
know whether it ever happens maliciously. We have shown that if it happens,
one profile sees it and another cannot.

\section{Conclusion}

An Agent Bill of Materials taken at task start is a constant function against
the substitution it most needs to catch. Making it a running commitment fixes
that for the capability surface: a tool list that changes between approval and
call moves a 32-byte register, and the event names the component that moved.
We measured what re-measurement costs locally --- around
\SI{0.7}{\micro\second} per tool, a hundredth of the per-action budget at
realistic list sizes --- and we did not measure the network round trip that
will decide whether it is affordable.

The same construction does nothing for content. A poisoned document that is
validly sourced and faithfully recorded moves the register exactly as a benign
one does, and the commitment records arrival rather than detecting anything.
We assert that negative in the test suite so it cannot quietly become a
positive, and we have said plainly that the harness re-states this result
rather than discovering it.

So the answer we would give P10 is narrower than the question hoped for:

\begin{quote}
A running composition commitment is a change-detection control over the
capability surface, and an inventory control over content.
\end{quote}

The second half of that sentence is the part worth carrying forward. It means
the memory- and context-poisoning rows of the threat model are not addressed
by any amount of additional rigour in the bill of materials, because the
problem there is not that we failed to record the component --- we recorded it
perfectly --- but that recording a component is not the same as knowing what
it did. That is a different research programme, and this tool's contribution
to it is to have declined to pretend otherwise, in a field on the profile
where the pretence would have gone.

\vspace{1em}
\noindent\textcolor{aaimuted}{\rule{\textwidth}{0.6pt}}
\vspace{0.4em}

\noindent{\footnotesize\textbf{Availability.} The implementation, the three
profiles, the detection harness and the benchmark are Apache-2.0 licensed
\cite{spectrumtool}. The detection outcomes in Table~\ref{tab:detect} are
reproduced by \texttt{spectrum detect}; Table~\ref{tab:cost} by
\texttt{spectrum bench-mcp}. Every test in the repository runs offline, with
no network, no credentials and no socket.}

\bibliographystyle{plain}
\bibliography{references}

\end{document}
